% Where a frame's time goes, for a 64 byte payload. Every bar is an estimate
% from the instruction counts. Nothing here has been measured yet, and the
% figure says so rather than letting an estimate read as a result.
\begin{tikzpicture}[font=\sffamily\small]
\node[font=\sffamily\bfseries\small, anchor=west] at (-20mm,52mm)
  {Cycles for one 64 byte payload at 280 MHz, estimated};

% ---- the grid goes down first so that nothing is drawn over a bar
\draw[taxis] (0mm,0mm) -- (126mm,0mm);
\foreach \x/\t in {0/0, 38.5/1000, 77/2000, 115.4/3000}{
  \draw[tick] (\x mm,0mm) -- (\x mm,48mm);
  \node[font=\tiny, anchor=north] at (\x mm,-1mm) {\t};}
\node[font=\sffamily\scriptsize, anchor=north west] at (0mm,-5mm) {cycles};

% 1 mm is 26 cycles
\foreach \y/\name/\w/\c/\s in {%
  42/byte stuffing/12.3/320/usrcol,
  33/CRC bitwise/119.2/3100/hwcol,
  24/CRC table/15.4/400/hwcol,
  15/{CRC unit, byte writes}/10.8/280/kercol,
  6/{CRC unit, word writes}/3.5/90/kercol}{
  \node[signame] at (-2mm,{\y mm+2.5mm}) {\name};
  \draw[signal, fill=\s] (0mm,\y mm) rectangle ({\w mm},{\y mm+5mm});
  \node[font=\sffamily\scriptsize, anchor=west] at ({\w mm+2mm},{\y mm+2.5mm})
    {\c};}

\node[font=\sffamily\scriptsize, anchor=west, text width=48mm, fill=white,
      inner sep=1.5pt] at (26mm,20.5mm)
  {plus 512 bytes of flash for the table, which is the other half of its
   speed};

\node[umlnote, text width=56mm, anchor=north west] at (-20mm,-10mm)
  {The interesting comparison is not hardware against software. It is the last
   two bars: a peripheral fed one byte at a time spends most of its advantage
   on bus accesses, and how much a wider write recovers is the question worth
   measuring.};
\node[umlnote, text width=56mm, anchor=north west] at (44mm,-10mm)
  {Every bar here is arithmetic from the instruction counts and none of it has
   been measured. The instrument that fills these in is the cycle counter,
   around a thousand repetitions, in step 9.};
\end{tikzpicture}
